\begin{examplebox}
	\textbf{\textcolor{CExample}{例 1（基础）}}
	\textbf{\textcolor{CExample}{题目：}}
	求函数 $f(x)=x^3-3x$ 的极值点，并判断是极大值还是极小值。
	\textbf{\textcolor{CExample}{解题步骤：}}
	\begin{description}[style=nextline, leftmargin=2.8em, labelsep=0.5em, itemsep=0.45em]
		\item[\textbf{第一步（求导与求根）：}]
			$f'(x)=3x^2-3=3(x-1)(x+1)$。令 $f'(x)=0$ 得 $x=-1,1$。
			\par\textbf{理由：} 极值点必要条件是导数为零。
		\item[\textbf{第二步（变号检验）：}]
			取 $x=-2$，$f'(-2)=9>0$；取 $x=0$，$f'(0)=-3<0$，故在 $x=-1$ 左右异号。同理，在 $x=1$ 左右：$x=0$ 时 $f'(0)=-3<0$，$x=2$ 时 $f'(2)=9>0$，异号。
			\par\textbf{理由：} 极值点充要条件是变号。
		\item[\textbf{第三步（结论与计数）：}]
			$x=-1$ 左正右负为极大值点，$x=1$ 左负右正为极小值点；两个驻点均变号，极值点个数为2。
			\par\textbf{理由：} 变号方向决定极大值和极小值。
	\end{description}
	\textbf{\textcolor{CExample}{关键结论：}}
	\textbf{极值点个数等于变号根个数。}

	\medskip
	\textbf{\textcolor{CExample}{例 2（稍变形）}}
	\textbf{\textcolor{CExample}{题目：}}
	判断函数 $f(x)=x^3$ 是否存在极值点，并说明理由。
	\textbf{\textcolor{CExample}{解题步骤：}}
	\begin{description}[style=nextline, leftmargin=2.8em, labelsep=0.5em, itemsep=0.45em]
		\item[\textbf{第一步（求导）：}]
			$f'(x)=3x^2$。
			\par\textbf{理由：} 幂函数求导。
		\item[\textbf{第二步（求根）：}]
			令 $f'(x)=0$ 得 $x=0$。
			\par\textbf{理由：} 极值必要条件检查。
		\item[\textbf{第三步（变号检验）：}]
			取 $x=-1$，$f'(-1)=3>0$；取 $x=1$，$f'(1)=3>0$；左右同号，故 $x=0$ 不是极值点。所以函数无极值点。
			\par\textbf{理由：} 根据充要条件，不变号就不是极值点。
	\end{description}
	\textbf{\textcolor{CExample}{关键结论：}}
	\textbf{驻点不一定都是极值点，必须验证变号。}

	\medskip
	\textbf{\textcolor{CExample}{例 3（综合应用）}}
	\textbf{\textcolor{CExample}{题目：}}
	已知函数 $f(x)=x^3-3ax$（$a\in\mathbb{R}$），讨论极值点的个数。
	\textbf{\textcolor{CExample}{解题步骤：}}
	\begin{description}[style=nextline, leftmargin=2.8em, labelsep=0.5em, itemsep=0.45em]
		\item[\textbf{第一步（求导）：}]
			$f'(x)=3x^2-3a=3(x^2-a)$。
			\par\textbf{理由：} 求导法则。
		\item[\textbf{第二步（根的讨论）：}]
			当 $a>0$ 时，$f'(x)=0$ 有两个实根 $\pm\sqrt{a}$；当 $a\le 0$ 时，$f'(x)=0$ 无实根（$a<0$）或一个根 $x=0$（$a=0$）。
			\par\textbf{理由：} 二次方程根的情况。
		\item[\textbf{第三步（变号检验与结论）：}]
			若 $a>0$，则 $x<-\sqrt{a}$ 时 $f'(x)>0$，$-\sqrt{a}<x<\sqrt{a}$ 时 $f'(x)<0$，$x>\sqrt{a}$ 时 $f'(x)>0$。故 $-\sqrt{a}$ 处左正右负（极大值点），$\sqrt{a}$ 处左负右正（极小值点），极值点个数为2。若 $a=0$，则 $f'(x)=3x^2\ge 0$，$x=0$ 左右 $f'(x)\ge 0$，不变号，不是极值点，个数为0。若 $a<0$，则 $f'(x)>0$ 恒成立，无极值点，个数为0。
			\par\textbf{理由：} 根据变号充要条件。
	\end{description}
	\textbf{\textcolor{CExample}{关键结论：}}
	\textbf{极值点个数 $=
	\begin{cases}
		2, & a>0,\\0, & a\le 0.
	\end{cases}
	$}
\end{examplebox}
